\section*{Pseudocódigo del Algoritmo Pan-Tompkins}

\subsection{1. Etapas de Filtrado y Procesamiento}
En esta primera etapa se desarrollaron las funciones de filtrado de datos y procesamiento de muestras para despues ser aplicadas con la ADT
\begin{algorithm}[H]
\caption{Procesamiento de Muestra}
\begin{algorithmic}[1]
\Function{ProcesarMuestra}{$m$}
    \If{not tieneOffset}
        \State offset $\leftarrow$ $m$
        \State tieneOffset $\leftarrow$ true
    \EndIf
    \State $m \leftarrow m - offset$
    \State pb $\leftarrow$ \Call{PasaBajos}{$m$}
    \State pa $\leftarrow$ \Call{PasaAltos}{pb}
    \State derivar $\leftarrow$ \Call{Derivada}{pa}
    \State elevarCuadrado $\leftarrow$ \Call{Cuadrado}{derivar}
    \State ventana $\leftarrow$ \Call{IntegracionVentana}{elevarCuadrado}
    \State \Return ventana
\EndFunction
\end{algorithmic}
\end{algorithm}

\begin{algorithm}[H]
\caption{Filtro Pasa-Bajos}
\begin{algorithmic}[1]
\Function{PasaBajos}{$x$}
    \State Desplazar y almacenar $x$ en buffer Pasa-Bajos
    \State $y \leftarrow 2 \cdot y[n-1] - y[n-2] + x[n] - 2 \cdot x[n-6] + x[n-12]$
    \State Almacenar $y$ en buffer de salidas Pasa-Bajos
    \State \Return $y$
\EndFunction
\end{algorithmic}
\end{algorithm}

\begin{algorithm}[H]
\caption{Filtro Pasa-Altos}
\begin{algorithmic}[1]
\Function{PasaAltos}{$x$}
    \State Desplazar y almacenar $x$ en buffer Pasa-Altos
    \State $y \leftarrow y[n-1] - \frac{x[n]}{32.0} + x[n-16] - x[n-15] + \frac{x[n-32]}{32.0}$
    \State Almacenar $y$ en buffer de salidas Pasa-Altos
    \State \Return $y$
\EndFunction
\end{algorithmic}
\end{algorithm}

\begin{algorithm}[H]
\caption{Derivada}
\begin{algorithmic}[1]
\Function{Derivada}{$x$}
    \State Desplazar y almacenar $x$ en buffer de Derivada
    \State $y \leftarrow 0.125 \cdot (2 \cdot x[n] + x[n-1] - x[n-3] - 2 \cdot x[n-4])$
    \State \Return $y$
\EndFunction
\end{algorithmic}
\end{algorithm}

\begin{algorithm}[H]
\caption{Cuadrado e Integración de Ventana}
\begin{algorithmic}[1]
\Function{Cuadrado}{$x$}
    \State \Return $x \cdot x$
\EndFunction

\Function{IntegracionVentana}{$x$}
    \State datoAntiguo $\leftarrow$ Sacar el dato más antiguo del buffer de integración
    \State suma\_int $\leftarrow$ suma\_int + $x$ - datoAntiguo
    \State Desplazar y almacenar $x$ en buffer de integración
    \State \Return suma\_int / window\_size
\EndFunction
\end{algorithmic}
\end{algorithm}

\subsection{2. Fase de Decisión y Lógica de Umbrales}
En esta parte del algoritmo a partir de los umbrales se decide el tipo de entrada y se clasifica, actualizando los umbrales a partir de la informacion obtenida por este proceso.
\begin{algorithm}[H]
\caption{Fase de Decisión}
\begin{algorithmic}[1]
\Procedure{Decidir}{$v$}
    \State pendMax $\leftarrow$ $\max(\text{pendMax}, v - prev1)$
    \If{not inicializado}
        \State maxInicializacion $\leftarrow$ $\max(\text{maxInicializacion}, v)$
        \State sumaInicializacion $\leftarrow$ sumaInicializacion + $v$
        \If{$n + 1 \geq \text{muestrasInicializacion}$}
            \State spki $\leftarrow$ $0.25 \cdot \text{maxInicializacion}$
            \State npki $\leftarrow$ $0.5 \cdot (\text{sumaInicializacion} / \text{muestrasInicializacion})$
            \State \Call{ActualizarUmbrales}{}
            \State pendMax $\leftarrow$ $0$
            \State inicializado $\leftarrow$ true
        \EndIf
    \Else
        \If{$prev1 > prev2$ and $prev1 \geq v$}
            \State \Call{Clasificar}{$prev1, n - 1$}
            \State pendMax $\leftarrow$ $0$
        \EndIf
        \State \Call{Searchback}{}
    \EndIf
    \State $prev2 \leftarrow prev1$
    \State $prev1 \leftarrow v$
    \State $n \leftarrow n + 1$
\EndProcedure

\Procedure{ActualizarUmbrales}{}
    \State thresholdI $\leftarrow$ npki + $0.25 \cdot (spki - npki)$
    \State thresholdII $\leftarrow$ $0.5 \cdot thresholdI$
\EndProcedure
\end{algorithmic}
\end{algorithm}

\subsection{3. Clasificación de Picos y Searchback}
Este ultimo bloque del algoritmo se encarga de clasificar el tipo de pico a base de su pentiente y nuevamente actualizar los umbrales a partir de la informacion obtenida
\begin{algorithm}[H]
\caption{Clasificación de Picos y Searchback}
\begin{algorithmic}[1]
\Procedure{Clasificar}{pico, idx}
    \State periodoRefactor $\leftarrow$ $\text{round}(0.200 \cdot fs)$
    \State ventanaT $\leftarrow$ $\text{round}(0.360 \cdot fs)$
    \State desde $\leftarrow$ idx - ultimoQRS

    \If{ultimoQRS $\geq 0$ and desde $<$ periodoRefactor}
        \State \textbf{return}
    \EndIf

    \If{pico $>$ thresholdI}
        \State ondaT $\leftarrow$ (ultimoQRS $\geq 0$) and (desde $<$ ventanaT) and (pendMax $< 0.5 \cdot pendQRS$)
        \If{ondaT}
            \State npki $\leftarrow$ $0.125 \cdot pico + 0.875 \cdot npki$
        \Else
            \State spki $\leftarrow$ $0.125 \cdot pico + 0.875 \cdot spki$
            \State \Call{RegistrarQRS}{idx, pendMax}
        \EndIf
    \Else
        \State npki $\leftarrow$ $0.125 \cdot pico + 0.875 \cdot npki$
        \If{pico $>$ thresholdII and pico $>$ candPico}
            \State candPico $\leftarrow$ pico
            \State candIdx $\leftarrow$ idx
            \State candPend $\leftarrow$ pendMax
        \EndIf
    \EndIf

    \State \Call{ActualizarUmbrales}{}
\EndProcedure

\Procedure{Searchback}{}
    \If{candIdx $\geq 0$ and haySuficientesIntervalosRR and $(n - ultimoQRS) > 1.66 \cdot rrReferencia()$}
        \State spki $\leftarrow$ $0.25 \cdot candPico + 0.75 \cdot spki$
        \State \Call{RegistrarQRS}{candIdx, candPend}
        \State \Call{ActualizarUmbrales}{}
    \EndIf
\EndProcedure
\end{algorithmic}
\end{algorithm}